\section{Error Decomposition}
\label{app:analysis}

This appendix derives the decomposition of Eq.~\mainref{eq:budget} that motivates
the design of \ours (Sec.~\mainref{sec:predict}). It describes how the errors of
an estimate scale with the lead and the motion; it is not a model of success.

\subsection{Translation}
\label{app:budget_translation}
Let the true target position at the execution time be
$\vect p(\tapp)=\vect p(\tau_t)+\vect v\Delta+\boldsymbol\rho$, where $\vect v$ is
the velocity at $\tau_t$ and $\boldsymbol\rho$ the remaining motion, and let the
policy's estimate be
$\hat{\vect p}(\tapp)=\hat{\vect p}(\tau_t)+\hat{\vect v}\hat\Delta+\hat{\boldsymbol\rho}$.
Subtracting and adding $\vect v\hat\Delta$ gives the identity
\begin{equation}
\begin{aligned}
\hat{\vect p}(\tapp)-\vect p(\tapp)
={}&\big[\hat{\vect p}(\tau_t)-\vect p(\tau_t)\big]+(\hat{\vect v}-\vect v)\hat\Delta\\
&+\vect v(\hat\Delta-\Delta)+(\hat{\boldsymbol\rho}-\boldsymbol\rho),
\end{aligned}
\label{eq:budget_identity}
\end{equation}
and the triangle inequality yields Eq.~\mainref{eq:budget} with
$e_v=\|\hat{\vect v}-\vect v\|$ and $e_{\rm dyn}=\|\hat{\boldsymbol\rho}-\boldsymbol\rho\|$.
The residual $\boldsymbol\rho$ is exactly the part of the motion that
constant-velocity extrapolation misses: under a constant acceleration $\vect a$,
$\boldsymbol\rho=\tfrac12\vect a\Delta^2$, so an acceleration estimate from the
history removes part of $e_{\rm dyn}$, whereas impacts, rolling contact, and friction
changes leave a remainder that only a model of the dynamics can reduce. For a
learned predictor, the split of its estimate into $\hat{\vect v}\hat\Delta$ and
$\hat{\boldsymbol\rho}$ is not unique, so the decomposition states which information
each term requires rather than how a particular network uses it.

\subsection{Rotation and the Contact Point}
\label{app:budget_rotation}
With $R(\tapp)=\Exp(\boldsymbol\omega\Delta+\boldsymbol\eta)R(\tau_t)$ and
$\hat R(\tapp)=\Exp(\hat{\boldsymbol\omega}\hat\Delta+\hat{\boldsymbol\eta})\hat R(\tau_t)$,
rotations about a common fixed axis compose additively, and
\begin{equation}
\theta=\big\|\Log(\hat R(\tapp)R(\tapp)^\top)^{\vee}\big\|
\le\theta_{\rm loc}+\hat\Delta\,e_\omega+\|\boldsymbol\omega\|\,e_\tau+e_\eta ,
\label{eq:budget_rotation}
\end{equation}
with the current orientation error $\theta_{\rm loc}$ and $e_\omega$ and $e_\eta$
defined as for translation; for general rotations, the Baker--Campbell--Hausdorff
formula adds terms of second order in the angles. A contact point at $\vect c$ in the
object frame, with $\|\vect c\|=r$, has the world position $\vect p+R\vect c$, so its
error is $(\hat{\vect p}-\vect p)+(\hat RR^\top-I)R\vect c$. Since the operator norm
of $\hat RR^\top-I$ for a rotation by $\theta$ is $2\sin(\theta/2)$,
\begin{equation}
e_{\rm c}\le\|\hat{\vect p}-\vect p\|+2r\sin(\theta/2)\le\|\hat{\vect p}-\vect p\|+r\theta .
\label{eq:contact_app}
\end{equation}
Orientation errors therefore enter the contact scaled by the lever arm, and a grasp
or placement that depends on orientation needs the full SE(3) state of the target.

\subsection{Context: Obstacles and Drivers}
\label{app:budget_context}
\paragraph{Obstacles.} Suppose that the executed path keeps every point of the robot and the carried object
at least $m$ away from the \emph{predicted} positions of the context points it
passes. A collision at a context point then requires the true position of that
point to differ from its prediction by at least $m$. Under this assumption, context
errors below the clearance do not cause collisions, which is why a coarse estimate
of the surroundings can suffice when $m\gg\varepsilon$. The assumption concerns the
executed path: \ours does not enforce a margin explicitly, so the effect of coarse
context supervision is tested empirically by the ablations
(Sec.~\mainref{sec:principles}). A flow sampled on a grid can also miss structures
thinner than its spacing.

\paragraph{Drivers.} A driver, such as a conveyor, a ramp, or a striker, matters
through the motion it imparts to the target. Its effect on the contact therefore
enters Eq.~\mainref{eq:budget} through the target's own terms, mostly through
$e_{\rm dyn}$ when the driver changes the target's motion, as at an impact. A policy
needs to know the driver's motion only well enough to anticipate when and how it will
act on the target; the resulting motion of the target is what the dedicated SE(3)
forecast supervises.

\subsection{Scope}
\label{app:budget_scope}
The decomposition addresses the approach to a rigid target whose motion is exogenous
until contact. It does not cover deformable or articulated targets, whose motion is
not determined by six numbers, or the motion after contact, which depends on the
action itself and would call for action-conditioned prediction. Scenes in which the
role of a body changes over time, e.g., sequential multi-object tasks, are also
outside its scope.
